%!TEX TS-program = xelatex

% HSE Beamer Theme
% by Danil Fedorovykh
% http://hse.ru/staff/df
%
% Version 2.0 (English)
% January 2022

%%% Set up the free HSE Sans font
%%% https://www.hse.ru/info/brandbook/#font

\documentclass[aspectratio=169]{beamer}

\newbool{russian}
\booltrue{russian} % Uncomment if in Russian
\usepackage{HSE-theme/beamerthemeHSE} % Load HSE theme
\usepackage[no-math]{fontspec}      % fonts loading
\usepackage{caption}
\usepackage{subfigure}
\usepackage{subcaption}
\usepackage{hyperref}
\usepackage[dvipsnames]{xcolor}
\usepackage{ragged2e}
\captionsetup[figure]{labelformat=empty}
\captionsetup[subfigure]{labelformat=empty}
\renewcommand*{\thesubfigure}{(\arabic{subfigure})}
\setsansfont{HSE Sans}
%\graphicspath{{./images/}}
\graphicspath{{/home/llyy/Yandex.Disk/personal/knowledge/general/algorithms_course/repo/algorithms_course/6_data_structures/lection/images}}


%%% Информация об авторе и выступлении
\title[Title]{Алгоритмы и структуры данных} 
\subtitle{Лекция 6. Структуры данных}
\author[Author's name]{Илья Сергеевич Бычков\\ \smallskip \scriptsize \url{ibychkov@hse.ru}}
\institute{НИУ ВШЭ - Нижний Новгород}
\date{\today}


\begin{document}

\frame[plain]{\titlepage}

%%%%%%%%%%%%%%%%%%%%%%%%%%%%%%%%%%%%%%%%%%%%%%%%%%%%%%%%%%%%%%%%%%%%%%%%%%%%%%%%%%%%%%%%%%%%%%%%%%
\begin{frame}[c]
%\frametitle{A first slide}

\begin{center}
\Huge Лекция 6.

\Huge Структуры данных
\end{center}

\end{frame}

%%%%%%%%%%%%%%%%%%%%%%%%%%%%%%%%%%%%%%%%%%%%%%%%%%%%%%%%%%%%%%%%%%%%%%%%%%%%%%%%%%%%%%%%%%%%%%%%%%
\begin{frame}
\frametitle{Структуры данных}
\framesubtitle{План лекции}

\begin{enumerate}
  \setcounter{enumi}{-1}
  \item{План лекции}
  \item{\textcolor{blue}{Стек с поддержкой минимума/максимума}}
  \item{Хэш-таблица (hash table)}
\end{enumerate}
\end{frame}


%%%%%%%%%%%%%%%%%%%%%%%%%%%%%%%%%%%%%%%%%%%%%%%%%%%%%%%%%%%%%%%%%%%%%%%%%%%%%%%%%%%%%%%%%%%%%%%%%%
\begin{frame}
\frametitle{Стек с поддержкой минимума}
\framesubtitle{Стек с поддержкой минимума}
\justifying

\textcolor{blue}{Задача:} Добавить возможность нахождения минимума в стеке за O (1), сохранив такой же асимптотику добавления и удаления элементов из стека.\newline\newline

\textcolor{blue}{Решение:} Хранить в стеке не сами элементы, а пары: элемент и минимум в стеке, начиная с этого элемента и ниже.


\end{frame}


%%%%%%%%%%%%%%%%%%%%%%%%%%%%%%%%%%%%%%%%%%%%%%%%%%%%%%%%%%%%%%%%%%%%%%%%%%%%%%%%%%%%%%%%%%%%%%%%%%
\begin{frame}
\frametitle{Структуры данных}
\framesubtitle{План лекции}

\begin{enumerate}
  \setcounter{enumi}{-1}
  \item{План лекции}
  \item{Стек с поддержкой минимума/максимума}
  \item{\textcolor{blue}{Хэш-таблица (hash table)}}
\end{enumerate}
\end{frame}


%%%%%%%%%%%%%%%%%%%%%%%%%%%%%%%%%%%%%%%%%%%%%%%%%%%%%%%%%%%%%%%%%%%%%%%%%%%%%%%%%%%%%%%%%%%%%%%%%%
\begin{frame}
\frametitle{Ассоциативный массив}
\framesubtitle{Ключи и значения}
\justifying
\textcolor{blue}{Задача}\newline
Хранить данные и уникальные идентификаторы, по которым эти данные можно различать. Идентификатор будем называть \textcolor{red}{ключ}, а данные - \textcolor{red}{значение}. Таким образом, необходимо уметь хранить пары (ключ, значение) и работать с ними. \newline

\textcolor{blue}{Абстракция ключ-значение}
\begin{itemize}
  \item{Вставка значения “по ключу”}
  \item{Дан ключ, необходимо найти значение}
\end{itemize}

\end{frame}

%%%%%%%%%%%%%%%%%%%%%%%%%%%%%%%%%%%%%%%%%%%%%%%%%%%%%%%%%%%%%%%%%%%%%%%%%%%%%%%%%%%%%%%%%%%%%%%%%%
\begin{frame}
\frametitle{Ассоциативный массив}
\framesubtitle{Direct-address tables}
\justifying
\textcolor{blue}{Direct-address tables}\newline
Простое решение в стиле сортировки подсчётом (counting sort). \newline Ключи - целые числа в некотором интервале [0, n - 1]. \newline Значения - произвольный объект. \newline

\textcolor{blue}{Абстракция ключ-значение}
\begin{itemize}
  \item{Вставка значения “по ключу”}
  \item{Дан ключ, необходимо найти значение}\newline
\end{itemize}

Операции поиска, вставки и удаления работают за O(1), так как они реализованы в виде работы с конкретными элементами массива.


\end{frame}

%%%%%%%%%%%%%%%%%%%%%%%%%%%%%%%%%%%%%%%%%%%%%%%%%%%%%%%%%%%%%%%%%%%%%%%%%%%%%%%%%%%%%%%%%%%%%%%%%%
\begin{frame}
\frametitle{Ассоциативный массив}
\framesubtitle{Direct-address tables}
\justifying
\textcolor{blue}{Direct-address tables}\newline

\begin{figure}
    \captionsetup[subfigure]{labelformat=empty}
    \centering
        \subfigure[{ \scriptsize Источник - \href{https://web.stanford.edu/class/archive/cs/cs161/cs161.1168/lecture9.pdf}{Stanford (Jessica Su) - Cormen et al.}}]{\includegraphics[width=100mm, height=50mm]{direct_address}} 
\end{figure}
\end{frame}

%%%%%%%%%%%%%%%%%%%%%%%%%%%%%%%%%%%%%%%%%%%%%%%%%%%%%%%%%%%%%%%%%%%%%%%%%%%%%%%%%%%%%%%%%%%%%%%%%%
\begin{frame}
\frametitle{Ассоциативный массив}
\framesubtitle{Direct-address tables}
\justifying
\textcolor{blue}{Direct-address tables}\newline

Операции поиска, вставки и удаления работают за O(1), так как они реализованы в виде работы с конкретными элементами массива.
\begin{figure}
    \captionsetup[subfigure]{labelformat=empty}
    \centering
        \subfigure[{ \scriptsize Источник - \href{https://web.stanford.edu/class/archive/cs/cs161/cs161.1168/lecture9.pdf}{Stanford (Jessica Su) - Cormen et al.}}]{\includegraphics[width=65mm, height=35mm]{direct_address_pseudo}} 
\end{figure}
\end{frame}

%%%%%%%%%%%%%%%%%%%%%%%%%%%%%%%%%%%%%%%%%%%%%%%%%%%%%%%%%%%%%%%%%%%%%%%%%%%%%%%%%%%%%%%%%%%%%%%%%%
\begin{frame}
\frametitle{Ассоциативный массив}
\framesubtitle{Ассоциативный массив}
\justifying
\textcolor{red}{Ассоциативный массив/словарь/отображение (dictionary/map)}\newline
это структура данных, обобщающая понятие массива, позволяющая хранить и \newline обрабатывать данные по произвольным \textcolor{red}{ключам}.\newline

\textcolor{blue}{Обобщение массива} Индексы массива не обязательно могут быть числами \newline от 0 до N-1\newline

Поддержка в современных языках
\begin{itemize}
\item{Сторонние библиотеки: C, VisualBasic, Standard ML, bash, ...}
\item{Встроенные библиотеки: Java, C\#, C++, Scala, ...}
\item{Поддержка на уровне языка: Awk, Perl, PHP, Tcl, JavaScript, Python, Ruby, Lua}
\end{itemize}

\end{frame}



%%%%%%%%%%%%%%%%%%%%%%%%%%%%%%%%%%%%%%%%%%%%%%%%%%%%%%%%%%%%%%%%%%%%%%%%%%%%%%%%%%%%%%%%%%%%%%%%%%
\begin{frame}
\frametitle{Ассоциативный массив}
\framesubtitle{Ассоциативный массив}

\textcolor{blue}{Пример - DNS lookup}
\begin{itemize}
  \item{Добавить URL с настоящим адресом}
  \item{Дан URL, необходимо узнать адрес}
\end{itemize}
\begin{figure}
    \captionsetup[subfigure]{labelformat=empty}
    \centering
        \subfigure[{ \scriptsize Источник - \href{https://algs4.cs.princeton.edu/home/}{Sedjwick, Wayne - algs4.cs.princeton.edu}}]{\includegraphics[width=70mm, height=20mm]{dict_dns_example}} 
\end{figure}
\end{frame}

%%%%%%%%%%%%%%%%%%%%%%%%%%%%%%%%%%%%%%%%%%%%%%%%%%%%%%%%%%%%%%%%%%%%%%%%%%%%%%%%%%%%%%%%%%%%%%%%%%
\begin{frame}
\frametitle{Ассоциативный массив}
\framesubtitle{Ассоциативный массив}
\textcolor{blue}{Ассоциативный массив.} Каждый элемент данных (значение) ассоциируется с некоторым идентификатором (ключ)\newline\newline
Функциональность:
\begin{itemize}
\item{put(key, value), обычно a[key] = value}
\item{put(key, value) перезаписывает старое значение новым}
\item{get(key), обычно a[key]}
\item{get(key) возвращает null если не существует значения по заданному ключу}\newline
\end{itemize}

Легкая реализация contains(key) и delete(key)

\end{frame}

%%%%%%%%%%%%%%%%%%%%%%%%%%%%%%%%%%%%%%%%%%%%%%%%%%%%%%%%%%%%%%%%%%%%%%%%%%%%%%%%%%%%%%%%%%%%%%%%%%
\begin{frame}
\frametitle{Ассоциативный массив}
\framesubtitle{Возможные реализации}

\textcolor{blue}{Структура данных.} Связный список пар ключ-значение\newline\newline
\textcolor{blue}{Поиск.} Исследуем все элементы пока не найдем совпадение или не закончится список\newline\newline
\textcolor{blue}{Вставка.} Исследуем все элементы пока не найдем совпадения, если не нашли – вставляем в начало\newline

\end{frame}

%%%%%%%%%%%%%%%%%%%%%%%%%%%%%%%%%%%%%%%%%%%%%%%%%%%%%%%%%%%%%%%%%%%%%%%%%%%%%%%%%%%%%%%%%%%%%%%%%%
\begin{frame}
\frametitle{Ассоциативный массив}
\framesubtitle{Возможные реализации}

\textcolor{blue}{Структура данных.} Связный список пар ключ-значение\newline
\begin{figure}
    \captionsetup[subfigure]{labelformat=empty}
    \centering
        \subfigure[{ \scriptsize Источник - \href{https://algs4.cs.princeton.edu/home/}{Sedjwick, Wayne - algs4.cs.princeton.edu}}]{\includegraphics[width=105mm, height=55mm]{dict_ll_impl}} 
\end{figure}
\end{frame}

%%%%%%%%%%%%%%%%%%%%%%%%%%%%%%%%%%%%%%%%%%%%%%%%%%%%%%%%%%%%%%%%%%%%%%%%%%%%%%%%%%%%%%%%%%%%%%%%%%
\begin{frame}
\frametitle{Ассоциативный массив}
\framesubtitle{Возможные реализации}

\textcolor{blue}{Структура данных.} Сохраняем элементы в виде таблицы(массива) при помощи \textcolor{red}{хэш-функции}.\newline\newline
\textcolor{blue}{Хэш-функция.} Функция для вычисления индекса массива по ключу.

\begin{figure}
    \captionsetup[subfigure]{labelformat=empty}
    \centering
        \subfigure[{ \scriptsize Источник - \href{https://algs4.cs.princeton.edu/home/}{algs4.cs.princeton.edu}}]{\includegraphics[width=65mm, height=40mm]{hash_example}} 
\end{figure}
\end{frame}

%%%%%%%%%%%%%%%%%%%%%%%%%%%%%%%%%%%%%%%%%%%%%%%%%%%%%%%%%%%%%%%%%%%%%%%%%%%%%%%%%%%%%%%%%%%%%%%%%%
\begin{frame}
\frametitle{Ассоциативный массив}
\framesubtitle{Хэш-функция}

\textcolor{blue}{Хэш-функция}. Функция для вычисления индекса массива по ключу.

Проблемы. 
\begin{itemize}
\item{Вычисление хэш-функции}
\item{Тест на равенство: Проверить что два ключа равны}
\item{Разрешение коллизий: Обработка ситуации когда два ключа имеют одинаковое значение хэш-функции}\newline
\end{itemize}

\textcolor{blue}{Classic space-time tradeoff}

\begin{itemize}
\item{Нет ограничений на память: тривиальная хэш-функция, где значение ключа и есть индекс}
\item{Нет ограничений на время: тривиальный алгоритм разрешения коллизий – использование линейного поиска}
\item{Ограничения: hashing (the real world)}\newline
\end{itemize}
\end{frame}

%%%%%%%%%%%%%%%%%%%%%%%%%%%%%%%%%%%%%%%%%%%%%%%%%%%%%%%%%%%%%%%%%%%%%%%%%%%%%%%%%%%%%%%%%%%%%%%%%%
\begin{frame}
\frametitle{Ассоциативный массив}
\framesubtitle{Хэш-функция}

\textcolor{blue}{Хэш-функция}\newline Идеальный случай: ключи равномерно распределены по массиву\newline


Эффективность по вычислениям\newline

\textcolor{blue}{Пример}. Телефонные номера\newline\newline
Плохой вариант: первые 3 цифры\newline
Хороший вариант: последние 3 цифры\newline

Нетривиальная задача на практике. \newline Необходим различный подход для разных типов ключей
\end{frame}

%%%%%%%%%%%%%%%%%%%%%%%%%%%%%%%%%%%%%%%%%%%%%%%%%%%%%%%%%%%%%%%%%%%%%%%%%%%%%%%%%%%%%%%%%%%%%%%%%%
\begin{frame}
\frametitle{Ассоциативный массив}
\framesubtitle{Хэш-функция}

\textcolor{blue}{Хэш-функция.} Функция для вычисления индекса массива по ключу.

\begin{figure}
    \captionsetup[subfigure]{labelformat=empty}
    \centering
        \subfigure[{ \scriptsize Источник - \href{https://algs4.cs.princeton.edu/home/}{Sedjwick, Wayne - algs4.cs.princeton.edu}}]{\includegraphics[width=100mm, height=55mm]{hash_example_java}} 
\end{figure}
\end{frame}

%%%%%%%%%%%%%%%%%%%%%%%%%%%%%%%%%%%%%%%%%%%%%%%%%%%%%%%%%%%%%%%%%%%%%%%%%%%%%%%%%%%%%%%%%%%%%%%%%%
\begin{frame}
\frametitle{Ассоциативный массив}
\framesubtitle{Хэш-функция}

\textcolor{blue}{Хэширование строк} \newline
\begin{itemize}
\item{A=a[0]a[1]a[2]…a[i] – символы}
\item{Использовать все символы – хорошая практика, иначе возможно большое количество коллизий}
\item{Например, если a[0] не используется, hash(“aa”) = hash(“ba”) = … = hash(“za”)}\newline
\end{itemize}


\textcolor{blue}{Полиномиальный хэш}

\begin{figure}
    \captionsetup[subfigure]{labelformat=empty}
    \centering
        \subfigure[{ \scriptsize Источник - Добавить}]{\includegraphics[width=120mm, height=9mm]{hash_poly}} 
\end{figure}
\end{frame}

%%%%%%%%%%%%%%%%%%%%%%%%%%%%%%%%%%%%%%%%%%%%%%%%%%%%%%%%%%%%%%%%%%%%%%%%%%%%%%%%%%%%%%%%%%%%%%%%%%
\begin{frame}
\frametitle{Ассоциативный массив}
\framesubtitle{Хэш-функция}

\textcolor{blue}{Хэш-функция.} Функция для вычисления индекса массива по ключу.

\begin{figure}
    \captionsetup[subfigure]{labelformat=empty}
    \centering
        \subfigure[{ \scriptsize Источник - \href{https://algs4.cs.princeton.edu/home/}{Sedjwick, Wayne - algs4.cs.princeton.edu}}]{\includegraphics[width=125mm, height=55mm]{hash_java_str}} 
\end{figure}
\end{frame}

%%%%%%%%%%%%%%%%%%%%%%%%%%%%%%%%%%%%%%%%%%%%%%%%%%%%%%%%%%%%%%%%%%%%%%%%%%%%%%%%%%%%%%%%%%%%%%%%%%
\begin{frame}
\frametitle{Ассоциативный массив}
\framesubtitle{Коллизия}

\textcolor{blue}{Коллизия.} Два различных ключа имеют одинаковое хэш-значение.\newline
\textcolor{blue}{Задача.} Придумать механизм для эффективного разрешения коллизий.

\begin{figure}
    \captionsetup[subfigure]{labelformat=empty}
    \centering
        \subfigure[{ \scriptsize Источник - Добавить}]{\includegraphics[width=85mm, height=55mm]{collision_names_example}} 
\end{figure}
\end{frame}

%%%%%%%%%%%%%%%%%%%%%%%%%%%%%%%%%%%%%%%%%%%%%%%%%%%%%%%%%%%%%%%%%%%%%%%%%%%%%%%%%%%%%%%%%%%%%%%%%%
\begin{frame}
\frametitle{Ассоциативный массив}
\framesubtitle{Метод цепочек}

\textcolor{blue}{Метод цепочек (chain hashing).} Используем массив из M < N связных списков. [H. P. Luhn, IBM 1953] .\newline\newline
\textcolor{blue}{Хэширование.} Отображение в целое число между 0 и M - 1. \newline\newline
\textcolor{blue}{Вставка.} вставка в начало цепи i (если элемента нет). \newline\newline
\textcolor{blue}{Поиск.} необходимость поиска только в цепи i. \newline\newline

\end{frame}


%%%%%%%%%%%%%%%%%%%%%%%%%%%%%%%%%%%%%%%%%%%%%%%%%%%%%%%%%%%%%%%%%%%%%%%%%%%%%%%%%%%%%%%%%%%%%%%%%%
\begin{frame}
\frametitle{Ассоциативный массив}
\framesubtitle{Метод цепочек}

\textcolor{blue}{Метод цепочек (chain hashing).} Используем массив из M < N связных списков. [H. P. Luhn, IBM 1953] .\newline

\begin{figure}
    \captionsetup[subfigure]{labelformat=empty}
    \centering
        \subfigure[{ \scriptsize Источник - \href{https://web.stanford.edu/class/archive/cs/cs161/cs161.1168/lecture9.pdf}{Stanford (Jessica Su) - Cormen et al.}}]{\includegraphics[width=110mm, height=50mm]{collision_chain_example2}}
         
\end{figure}
\end{frame}

%%%%%%%%%%%%%%%%%%%%%%%%%%%%%%%%%%%%%%%%%%%%%%%%%%%%%%%%%%%%%%%%%%%%%%%%%%%%%%%%%%%%%%%%%%%%%%%%%%
\begin{frame}
\frametitle{Ассоциативный массив}
\framesubtitle{Метод цепочек}

\textcolor{blue}{Метод цепочек (chain hashing).} Используем массив из M < N связных списков. [H. P. Luhn, IBM 1953] .\newline

\begin{figure}
    \captionsetup[subfigure]{labelformat=empty}
    \centering
        \subfigure[{ \scriptsize Источник - \href{https://algs4.cs.princeton.edu/home/}{Sedjwick, Wayne - algs4.cs.princeton.edu}}]{\includegraphics[width=90mm, height=50mm]{collision_chain_example1}} 
\end{figure}
\end{frame}

%%%%%%%%%%%%%%%%%%%%%%%%%%%%%%%%%%%%%%%%%%%%%%%%%%%%%%%%%%%%%%%%%%%%%%%%%%%%%%%%%%%%%%%%%%%%%%%%%%
\begin{frame}
\frametitle{Хэш-таблица}
\framesubtitle{Метод цепочек}

\textcolor{blue}{Издержки.} Число исследований (\textcolor{red}{probes}) для поиска и вставки пропорционально N / M. 

\begin{itemize}
\item{Слишком большое M ⇒ слишком много пустых цепочек}
\item{Слишком маленькое M ⇒ слишком длинные цепочки}
\item{Практический результат: M ~ N / 4 ⇒ константное время операций}\newline
\end{itemize}

\textcolor{blue}{Цель}. Средняя длинна цепи N / M - константа.
\begin{itemize}
\item{Удваиваем значение M когда N / M ≥ 8}
\item{Делим пополам значение M когда N / M ≤ 2}
\item{Повторное хэширование после изменения размеров}\newline
\end{itemize}

\end{frame}

%%%%%%%%%%%%%%%%%%%%%%%%%%%%%%%%%%%%%%%%%%%%%%%%%%%%%%%%%%%%%%%%%%%%%%%%%%%%%%%%%%%%%%%%%%%%%%%%%%
\begin{frame}
\frametitle{Хэш-таблица}
\framesubtitle{Метод цепочек}

\textcolor{blue}{Метод цепочек (chain hashing).} Используем массив из M < N связных списков. [H. P. Luhn, IBM 1953] .\newline

\begin{figure}
    \captionsetup[subfigure]{labelformat=empty}
    \centering
        \subfigure[{ \scriptsize Источник - \href{https://algs4.cs.princeton.edu/home/}{Sedjwick, Wayne - algs4.cs.princeton.edu}}]{\includegraphics[width=90mm, height=50mm]{collision_chain_resize}} 
\end{figure}
\end{frame}

%%%%%%%%%%%%%%%%%%%%%%%%%%%%%%%%%%%%%%%%%%%%%%%%%%%%%%%%%%%%%%%%%%%%%%%%%%%%%%%%%%%%%%%%%%%%%%%%%%
\begin{frame}
\frametitle{Хэш-таблица}
\framesubtitle{Метод цепочек}

Удаление элемента\newline

\begin{figure}
    \captionsetup[subfigure]{labelformat=empty}
    \centering
        \subfigure[{ \scriptsize Источник - \href{https://algs4.cs.princeton.edu/home/}{Sedjwick, Wayne - algs4.cs.princeton.edu}}]{\includegraphics[width=105mm, height=50mm]{collision_chain_remove}} 
\end{figure}
\end{frame}

%%%%%%%%%%%%%%%%%%%%%%%%%%%%%%%%%%%%%%%%%%%%%%%%%%%%%%%%%%%%%%%%%%%%%%%%%%%%%%%%%%%%%%%%%%%%%%%%%%
\begin{frame}
\frametitle{Хэш-таблица}
\framesubtitle{Метод открытой адресации}

\textcolor{blue}{Метод открытой адресации (open addressing).} При возникновении коллизии ищем первый доступный пустой слот и помещаем туда ключ/значение [Amdahl-Boehme-Rocherster-Samuel, IBM 1953] .

\begin{figure}
    \captionsetup[subfigure]{labelformat=empty}
    \centering
        \subfigure[{ \scriptsize Источник - \href{https://algs4.cs.princeton.edu/home/}{Sedjwick, Wayne - algs4.cs.princeton.edu}}]{\includegraphics[width=45mm, height=45mm]{open_addressing_table}} 
\end{figure}
\end{frame}

%%%%%%%%%%%%%%%%%%%%%%%%%%%%%%%%%%%%%%%%%%%%%%%%%%%%%%%%%%%%%%%%%%%%%%%%%%%%%%%%%%%%%%%%%%%%%%%%%%
\begin{frame}
\frametitle{Хэш-таблица}
\framesubtitle{Метод открытой адресации}

\textcolor{blue}{Метод открытой адресации (open addressing)}\newline

\textcolor{blue}{Хэширование.} Отображение ключей на интервал [0, М-1]\newline
\textcolor{blue}{Вставка.} Если индекс в массиве свободен – записываем значение туда. Иначе пробуем i+1, i+2 и тд\newline
\textcolor{blue}{Note.} Размер массива должен вместить все пары ключ-значение.\newline
search K\newline
hash(K) = 5

\begin{figure}
    \captionsetup[subfigure]{labelformat=empty}
    \centering
        \subfigure[{ \scriptsize Источник - \href{https://algs4.cs.princeton.edu/home/}{Sedjwick, Wayne - algs4.cs.princeton.edu}}]{\includegraphics[width=125mm, height=25mm]{open_addressing_example}} 
\end{figure}
\end{frame}

%%%%%%%%%%%%%%%%%%%%%%%%%%%%%%%%%%%%%%%%%%%%%%%%%%%%%%%%%%%%%%%%%%%%%%%%%%%%%%%%%%%%%%%%%%%%%%%%%%
\begin{frame}
\frametitle{Хэш-таблица}
\framesubtitle{Метод открытой адресации}

\begin{figure}
    \captionsetup[subfigure]{labelformat=empty}
    \centering
        \subfigure[{ \scriptsize Источник - \href{https://algs4.cs.princeton.edu/home/}{Sedjwick, Wayne - algs4.cs.princeton.edu}}]{\includegraphics[width=90mm, height=60mm]{open_addressing_linear_probing}} 
\end{figure}
\end{frame}

%%%%%%%%%%%%%%%%%%%%%%%%%%%%%%%%%%%%%%%%%%%%%%%%%%%%%%%%%%%%%%%%%%%%%%%%%%%%%%%%%%%%%%%%%%%%%%%%%%
\begin{frame}
\frametitle{Хэш-таблица}
\framesubtitle{Linear robing}

\begin{figure}
    \captionsetup[subfigure]{labelformat=empty}
    \centering
        \subfigure[{ \scriptsize Источник - \href{https://web.stanford.edu/class/archive/cs/cs161/cs161.1168/lecture9.pdf}{Stanford (Jessica Su) - Cormen et al.}}]{\includegraphics[width=140mm, height=20mm]{linear_probing}} 
\end{figure}

\end{frame}

%%%%%%%%%%%%%%%%%%%%%%%%%%%%%%%%%%%%%%%%%%%%%%%%%%%%%%%%%%%%%%%%%%%%%%%%%%%%%%%%%%%%%%%%%%%%%%%%%%
\begin{frame}
\frametitle{Хэш-таблица}
\framesubtitle{Quadratic robing}

\begin{figure}
    \captionsetup[subfigure]{labelformat=empty}
    \centering
        \subfigure[{ \scriptsize Источник - \href{https://web.stanford.edu/class/archive/cs/cs161/cs161.1168/lecture9.pdf}{Stanford (Jessica Su) - Cormen et al.}}]{\includegraphics[width=140mm, height=20mm]{quadratic_probing}} 
\end{figure}

\end{frame}

%%%%%%%%%%%%%%%%%%%%%%%%%%%%%%%%%%%%%%%%%%%%%%%%%%%%%%%%%%%%%%%%%%%%%%%%%%%%%%%%%%%%%%%%%%%%%%%%%%
\begin{frame}
\frametitle{Хэш-таблица}
\framesubtitle{Double hashing}

\begin{figure}
    \captionsetup[subfigure]{labelformat=empty}
    \centering
        \subfigure[{ \scriptsize Источник - \href{https://web.stanford.edu/class/archive/cs/cs161/cs161.1168/lecture9.pdf}{Stanford (Jessica Su) - Cormen et al.}}]{\includegraphics[width=140mm, height=15mm]{double_hashing}} 
\end{figure}

\end{frame}

%%%%%%%%%%%%%%%%%%%%%%%%%%%%%%%%%%%%%%%%%%%%%%%%%%%%%%%%%%%%%%%%%%%%%%%%%%%%%%%%%%%%%%%%%%%%%%%%%%
\begin{frame}
\frametitle{Хэш-таблица}
\framesubtitle{Метод открытой адресации}


Many improved versions have been studied. \newline

\textcolor{blue}{Two-probe hashing} [ separate-chaining variant ] \newline
Hash to two positions, insert key in shorter of the two chains. \newline

\textcolor{blue}{Double hashing} [ linear-probing variant ] \newline
Use linear probing, but skip a variable amount, not just 1 each time. \newline
Effectively eliminates clustering. \newline

\textcolor{blue}{Cuckoo hashing} [ linear-probing variant ] \newline
Hash key to two positions; insert key into either position; if occupied, reinsert displaced key into its alternative position. 


\end{frame}

\end{document}